\documentclass[../main.tex]{subfiles}
\graphicspath{{\subfix{../image/}}}

\begin{document}

\chapter{定向}\label{chapter:ch15-orientations}

本章的目的是介绍光滑流形一个微妙而重要的性质，称为定向。这个词源于拉丁文 oriens（“东方”），最初意为“朝东”，更一般地说意为“相对于周围环境的方位”。因此，一条直线或曲线的定向，只不过是沿着它选取一个方向。正如我们在第~11 章中所见，线积分的符号取决于沿曲线选取的优选方向。

数学家们把“定向”一词的含义推广到了更高维流形，即物体相对于其周围环境所能采取的两种不等同位置方式之间的选择。对二维流形而言，定向本质上就是选择哪个旋转方向应被视为“顺时针”、哪个为“逆时针”。对三维流形而言，它是在“左手性”与“右手性”之间的选择。定向的一般定义就是把这些日常概念推广到任意维数。

正如我们将在本章中看到的，一个向量空间恰有两种定向选择。在 $\mathbb{R}^n$ 中，存在一个我们都能认同的标准定向；但在其他向量空间中，必须作出一个任意选择。对流形而言，情况要复杂得多。在球面上，通过从外部观察表面，可以毫不含糊地断定哪个旋转方向是逆时针（\cref{figure:ch15-sphere}）。另一方面，Möbius 带（\cref{figure:ch15-mobius}）具有这样一个奇特的性质：在其表面上移动的图形可以回到起点，却变成了它的镜像，因而无法始终如一地断定该表面上两个可能的旋转方向中哪个叫“顺时针”、哪个叫“逆时针”，也无法断定哪一面是“正面”、哪一面是“背面”。在三维流形中，类似的现象将是一个惯用右手的人作了一次长途旅行后回来变成了左撇子。像球面这样可以选取一致定向的流形称为\emph{可定向的}；像 Möbius 带这样不能选取一致定向的流形称为\emph{不可定向的}。

在本章中，我们发展光滑流形的定向理论。它们有众多应用，最显著的是我们将在第~16 章中研究的流形上积分理论。

\subfile{Chapter-15/section-01.tex}
\subfile{Chapter-15/section-02.tex}
\subfile{Chapter-15/section-03.tex}
\subfile{Chapter-15/section-04.tex}
\subfile{Chapter-15/section-05.tex}

\end{document}
